\chapter{Measurement Defects}
\label{app:defects}

\keybox{what this appendix is for}{%
Eight defects were found in this project's measurement, by checks built for the
purpose, and each one produced numbers that looked entirely reasonable while
being wrong. They live in one chapter, away from Chapter~\ref{ch:results}, because a reader
following the argument wants the argument, and a reader auditing the work wants
all the forensics in one place.

Each entry follows the same shape: what the number said, what was actually
happening, how it was found, and what now prevents it. The last of those is the
point. A defect that is fixed but not covered by a check is a defect that will
recur under a different name --- and two of the eight below are the same mistake
at different denominators.}

\section*{The Eight, In One Table}

\begin{center}\small
\begin{tabularx}{\linewidth}{L{1}L{1}L{1}}
  \toprule
  what looked right & what was wrong & what catches it now \\
  \midrule
  The 3DGS arm reproduced 3DGS's published full-resolution figure &
  It still carried Mip-Splatting's screen-space opacity compensation, so it was
  anti-aliasing better than 3DGS at reduced scales &
  The switch is asserted on the \emph{compiled} rasteriser, not the source \\
  \addlinespace
  The opacity-compensation fix switch isolated one term &
  It removed three: the compensation, and a degeneracy guard in each of the
  forward and backward kernels &
  The guard is hoisted out of the branch in both kernels \\
  \addlinespace
  Conditioning and rendered quality agreed per protocol &
  The instrument and the trainer were cloning different branches, so the same
  protocol name meant two different captures &
  Both pin a branch; \texttt{test\_camera\_subset.py} asserts what each protocol
  selects \\
  \addlinespace
  The estimation term exceeded the floor everywhere &
  The constant $0.2$ was carried on one side of the comparison and not the other,
  inflating it fivefold &
  Proposition~\ref{prop:constant} is a unit test, passing to $1.6\times10^{-16}$ \\
  \addlinespace
  Table~\ref{tab:mtmt}'s 3DGS column &
  Every one of its $32$ cells was the mean of two different rasterisers &
  The generator refuses any table whose cells span more than one build \\
  \addlinespace
  Mip-Splatting reproduced to $0.50$\,dB --- L2 &
  A seven-scene mean compared against an eight-scene published figure. The true
  figure is $\stmtArmAMaxDelta$\,dB --- L1 &
  Against-published uses the arm's own scene set; between-arms uses the matched
  set; a test asserts all three outputs agree \\
  \addlinespace
  The offset was explained by GOF densification &
  Asserted for months, never tested. Measured, the effect is
  $\gofMaxEffect$\,dB and not consistently signed &
  The pre-GOF commit is a probe in the record (Table~\ref{tab:pregof}) \\
  \addlinespace
  The angular mask lost heavily to magnitude pruning &
  Against a \emph{free} control with more freedom than the angular mask has. At matched
  structure the margin is $\btwoBlockedGapMin$--$\btwoBlockedGapMax$\,dB &
  The control carries the angular mask's own nested-degree structure \\
  \bottomrule
\end{tabularx}
\end{center}

\medskip
Five of the eight are given in full below. The other three --- the denominator,
the densification explanation and the angular mask's unfair control --- cannot be separated
from the results they correct, so they are told where those results are:
\S\ref{sec:table1}, \S\ref{sec:densification} and \S\ref{sec:btwo}
respectively.

\paragraph{One Pattern, Six Times.} Six of the eight are the same error: two
quantities were compared that were not measured over the same thing. Two
rasterisers averaged into one cell. A seven-scene mean against an eight-scene
figure. A protocol name meaning two captures. A constant present on one side and
not the other. A pruning control with more freedom than the method it judged. In
every case the arithmetic was right and the comparison was not, and in every case
the result looked plausible --- which is why none of them was caught by reading
the numbers, and all of them were caught by a check that asked what the numbers
were means \emph{over}.

\section{The Fourth Check: The Rasteriser On The 3DGS Arm}

The single-codebase construction of \S\ref{sec:onecodebase} gives both arms
Mip-Splatting's rasteriser, and sets the 3DGS arm apart by zeroing \texttt{filter\_3D}
and taking \texttt{kernel\_size}~$=0.3$. That reproduces 3DGS's \emph{3D}
band-limit exactly --- which is what check~2 verifies --- and says nothing about
the \emph{2D} one. Mip-Splatting's screen-space filter dilates the projected
covariance by $k$ and then compensates the opacity by
\begin{equation}
  \rho \;=\; \sqrt{\frac{\lvert\Sigma'\rvert}{\lvert\Sigma' + kI\rvert}} ,
\end{equation}
so that dilation does not brighten the primitive. Vanilla 3DGS applies the same
fixed $0.3$ dilation and \textbf{no compensation at all}. Setting
\texttt{kernel\_size}~$=0.3$ reproduces the dilation magnitude and leaves $\rho$
in place, so the 3DGS arm kept an anti-aliasing term that 3DGS does not have --- and
the term it kept is precisely the one that protects a dilated primitive from
losing energy as the test sampling rate falls. An arm that anti-aliases better
than 3DGS degrades less than 3DGS, which is what the scale-collapse gate measured.

\paragraph{The Test That Separates The Two Explanations.} $\rho$ is
computed in CUDA, so the switch is a rasteriser change: a
\texttt{mip\_compensation} flag threaded through
\texttt{Gaussian\allowbreak Rasterization\allowbreak Settings} into
\texttt{computeCov2D}, which skips
the $\rho$ term while keeping the dilation. Everything else in that function is
identical to \texttt{graphdeco-inria/\allowbreak diff-gaussian-\allowbreak rasterization} line
for line.
The build asserts the compiled binary carries the field before any training
starts, because checking the source would prove nothing about the binary.

One scene, \texttt{lego}, 30\,000 iterations, everything else held fixed:

\input{generated/table-c1}

Full resolution barely moves: $\cOneBaseFull$ against $\cOneFull$, which is
what must happen, because at the training sampling rate the dilation is
negligible against the primitive's own extent and $\rho \approx 1$. Every
reduced scale collapses, and lands on the shape of the published 3DGS row:
$\cOneHalf$ against $26.95$, $\cOneQuarter$ against $21.38$, $\cOneEighth$
against $17.69$. The compensation was the whole of the defect.

\keybox{the densification hypothesis}{%
Both columns of Table~\ref{tab:c1} ran GOF's densification. If densification were
the cause of the shortfall, the right-hand column could not have collapsed --- and
it collapses by $\cOneCollapse$\,dB at $\tfrac18$, landing within $0.3$\,dB of
the published figure at every reduced scale. The earlier diagnosis is therefore
wrong, and it is left standing in this chapter because how a wrong diagnosis
survived matters: it was never given an experiment that could refute
it. It explained the anomaly as a property of the codebase, which is the most
comfortable kind of explanation and the one that most deserves a test.

One point is easy to read the wrong way. The comparison that settles this runs
\emph{between the two columns}; either column set against the published row
settles nothing. Both columns are \texttt{lego} and
the published row is an eight-scene mean, so the agreement at reduced scales is
corroboration and not a measurement. The measurement is that switching one CUDA
term moved $\tfrac18$ by $\cOneCollapse$\,dB and full resolution by less than
$0.05$.}


\section{What The Vanilla Arm Cost In Robustness}
\label{sec:shipfails}

Removing the compensation also removed a guard. The same CUDA branch that
computes $\rho$ carries \texttt{coef = 0} for a primitive whose projected
covariance is singular to the tolerance the determinants are clamped at --- a
numerical guard, not part of the anti-aliasing model. With the compensation
switched off, degenerate primitives survived at full opacity with an unbounded
screen-space radius, and two scenes died reproducibly, on both GPUs, with a CUDA
illegal memory access around iteration $3\,000$. Six trained normally, which is
what made it look transient.

Hoisting the guard out of the branch fixed one of the two. The other then failed
in the \emph{backward} pass instead, because the backward kernel still had its
copy of the guard inside the branch: the forward was contributing nothing for a
degenerate primitive while a gradient still flowed to its opacity. Zeroing a
forward contribution and leaving its gradient live is not a consistent pair, and
matching them is what the opacity-compensation fix switch should have done from the start --- it was
supposed to isolate one term and it removed three.

\paragraph{\texttt{ship} Does Not Complete, And Is Reported As Such.} With both
guards matched it fails a third time, and a fourth under instrumentation. The
first three tracebacks all pointed at \texttt{ssim}, in the \emph{next}
iteration's loss, which is a red herring: a CUDA fault is reported
asynchronously, at whichever call next touches the device. Re-run with
\texttt{CUDA\_LAUNCH\_BLOCKING=1} so that each launch is synchronous, it
localises instead to the forward rasteriser --- the fault surfaces on the first
tensor derived from its output, \texttt{radii > 0} in
\texttt{gaussian\_renderer}, with the backward pass never reached.

That clears three explanations. The backward kernel, which the second failure
had suggested, is clear: the fault surfaces before the backward pass is reached.
Memory exhaustion is clear: every attempt raises an illegal address and none
raises an allocation error, and \texttt{ship} completes on the Mip-Splatting arm
on the same $15\,360$\,MiB card at a measured peak of $13\,755$\,MiB. And the
onset moves between attempts, which clears the deterministic fault at a fixed
point this section claimed before it was instrumented.

The scene has now been attempted nine times on the vanilla arm --- four
diagnostic runs and then five consecutive attempts in one session, since a fault
whose onset moves is a fresh draw each time, so retrying it is a measurement. \textbf{All nine failed}, every one with
\texttt{cudaErrorIllegalAddress} and none with an allocation error, at
iterations
\[
  3\,000,\quad 7\,220,\quad 8\,580,\quad 10\,100,\quad 10\,210,\quad
  10\,320,\quad 11\,820,\quad 12\,100
\]
and one earlier still. What varies is \emph{when}, not \emph{whether}: the
failure is certain and its timing is not. That settles the question retrying was
asked to settle --- there is no lucky run to be had --- and it is why the scene
is reported as failed and left there.

What remains is an out-of-bounds access in the forward raster whose onset
depends on the optimisation trajectory, on much the largest model in the suite:
\texttt{ship} reaches $442\,235$ primitives on the Mip-Splatting arm, more than any other
Blender scene, at $\sim\!90\,\%$ of the card.

It is excluded from the single-scale 3DGS-arm rows, and because every
between-arm
table in this chapter is restricted to the scenes both arms have, it is excluded
from the Mip-Splatting arm's there too --- a seven-scene mean set against an eight-scene one compares two different things. Nothing was reduced to make it fit: the protocol's rule for a scene
that will not run is to record it and move on, and that rule does not become
optional when the scene is one of eight, as it would be one of nine.

Two things keep the exclusion narrow. \texttt{ship}
trains without incident on the \emph{multi-scale} protocol, on all four scales,
so Table~\ref{tab:mtmt} is a full eight-scene comparison and the failure is
specific to single-scale training. And it trains without incident on the Mip-Splatting arm, so
the fragility belongs to the vanilla band-limit semantics, with the scene, the
data and the harness all cleared. Both point the same way: removing the
compensation costs robustness on the one scene that sits closest to the limits of
the card, and it does so through the forward raster.

The mechanism is \emph{not} settled beyond that, and the earlier draft of this
section overstated it. Hoisting the degeneracy guard out of the compensation
branch was a real defect and a real fix --- it is what recovered \texttt{chair}
--- but it cannot also be the explanation for the failures that persist after
it, because the guard now fires identically whether the compensation is on or
off. What can be said from measurement is the conjunction above: forward raster,
varying onset, no allocation error, largest model in the suite, vanilla arm
only. Naming a cause more specific than that would be inference presented as
result, which is the thing \S\ref{sec:g2} already cost this project once.

\keybox{what makes a switch a control}{%
Each of those three removals was invisible: no assertion covered them, the
numbers that came back were plausible, and six of eight scenes were unaffected.
What exposed them was a failure that could not be waved through --- a crash --- and
then a failure that \emph{moved} from the forward pass to the backward one when
half the fix was applied, which is what identified the second half. A switch
meant to isolate a single term has to be checked against what it leaves alone,
not only against what it changes.}


\section{A Column Averaged Over Two Rasterisers}
\label{sec:blend}

The 3DGS arm's column here was re-measured after the opacity-compensation fix, and for a while afterwards each of
its $32$ cells averaged the old measurement together with the new one, where the
new one alone belonged there. The record is append-only by design (Table~\ref{tab:levels}): a
re-measured configuration leaves its predecessor in the file, marked, and the
tables skip anything marked. The marking rule identified the rows to retire by
the rung name in their notes --- and the re-run was the same rung, so its rows
carried the same name, the rule exempted both, and nothing was retired. Both
rows then satisfied the table's selector and were averaged.

Nothing about the output looked wrong. The two builds differ only in whether the
degeneracy guards in the forward and backward passes match, which moves a scene
by hundredths of a decibel on a capture this well conditioned, so the blended
column sat exactly where a plausible column sits. The same defect had quietly
taken hold in the anisotropic filter parity rows, where the two builds differed in whether the
filter floor was Mip-Splatting's own value or a re-derivation of it.

What was missing was not a better rule but any check at all on the invariant the
rule exists to maintain. There is one now: the generator refuses to emit a table
if any reported cell draws on more than one build, and it names the cell. It is
four lines, it would have caught both instances the day they appeared, and it
subsumes every future variant of the same mistake --- which a third special case
in the marking rule would not have.

\keybox{what append-only applies to}{%
Keeping superseded rows is what makes a result auditable, and this project has
been strict about it. But a row that is kept and not marked is a row that is
\emph{used}, and the marking is the only thing standing between an honest
archive and a silently pooled average. A discipline that depends on a marking
step needs the pooled average checked, not the marking step trusted.}


\section{Two Kernels, One Set Of Protocol Names}
\label{sec:branchdefect}

The instrument results and the PSNR results of \S\ref{sec:table3} were, for a
time, measurements of \emph{different captures} reported under the same protocol
names, and the mechanism matters, because nothing in either run could have
detected it.

The instrument kernel cloned the repository's default branch. The training kernel
clones the method branch. Between them, the definition of \texttt{arc} and
\texttt{cone} had changed: the default branch sizes them by angular \emph{width},
giving $20$ cameras across $81^\circ$ and $3$ across $30^\circ$; the method
branch sizes them by camera \emph{count}, giving $25$ across $93^\circ$ and $10$
across $54^\circ$. The change was deliberate and its stated reason is sound --- a
$20^\circ$ wedge catches three cameras on a Blender capture, and three views turns a conditioning experiment into a sparse-reconstruction
one. What was missed
is that only one of the two kernels was updated.

Each run was internally consistent, each printed the span it had built, and each
was correct about its own capture. The error was in reading them side by side,
and it survived because both numbers looked reasonable: a cone \emph{is} narrower
than an arc under either definition, so the ordering the chapter argues for held
either way.

\keybox{the cost of this defect}{%
It cost the first reading of the floor-binding and anisotropy gates --- and it bought the \texttt{pencil}
protocol. Rather than choose between the two definitions, both are kept: the
count-based \texttt{cone} at $51^\circ$, which is a trainable subset, and the
width-based \texttt{pencil} at $20^\circ$, which is the regime the derivation is
about. The crossover between them is where the estimation term overtakes the
floor, and with both in the table that crossover is measured directly. The instrument kernel is now pinned to the branch the training
kernel clones, and every protocol table carries the span the session that
produced the row actually built.}

\section{A Constant That Was Carried On One Side Only}
\label{sec:m2defect}

One more defect has to be stated before the numbers, because it changed them.
The estimation term is $(s\,\sigma_{\mathrm{pix}})^2/\lambda_i$, and
Proposition~\ref{prop:constant} fixes $(s\,\sigma_{\mathrm{pix}})^2 = 0.2$ ---
the same constant Mip-Splatting's floor $f_k^2 = 0.2\,(d_{\min}/f_{\max})^2$ is
built from. Both sides of the floor-occupancy instrument comparison therefore carry it. The instrument
carried it on the floor and dropped it from the estimation term, which inflated
that side by a factor of five, or $\sqrt{5} = 2.24$ in $\sigma$. The filter's own
implementation had always applied it.

The symptom was two measurements of one quantity disagreeing --- on the arc and
the cone, and not at the extremes, where a factor of five changes no verdict.
Nothing raised, and both numbers were plausible on their own.

\keybox{Proposition~1 as a test}{%
The identity that fixes the constant is checkable: at one view, on the optical
axis, the estimation term must \emph{equal} the floor. It now does, to
$1.6\times10^{-16}$, and \texttt{tools/test\_instruments.py} asserts it on every
run. Without the $0.2$ it is exactly five times the floor, so the test cannot
pass either way by accident. This is the over-determination the audit's M-2
named --- ``$s$ defaults to 1'' and ``$s\,\sigma_{\mathrm{pix}} = \sqrt{0.2}$''
cannot both be free --- surfacing as a live numerical error, the most useful form it could take; as an
ambiguity in the prose it would have stood for years.}


\section{Three Bugs The Parity Requirement Caught}
\label{sec:b1bugs}

The first implementation of the anisotropic filter reported exactly what
\S\ref{sec:instrumentresults} predicts internally --- the floor binding on
$100\,\%$ of primitives, the filter isotropic to two decimals --- and still
differed from Mip-Splatting by up to $1.22$\,dB, thirty times the seed spread of
\S\ref{sec:spread}. Section~\ref{ch:discussion} names that case in advance: a
loss on the full-orbit row is a bug, not a result, because
Proposition~\ref{prop:reduction} forbids it. Diagnosing it took three passes, and each step ruled something out.

\begin{enumerate}[leftmargin=1.4em,itemsep=3pt]
  \item \textbf{The floor.} The anisotropic filter re-derives $f_k$ from the camera geometry, so the
        first suspect was that its floor was not Mip-Splatting's. A cross-check
        was added --- \texttt{compute\_3D\_filter} runs alongside and the two are
        compared on every call --- and they agree to $1.16\times10^{-10}$
        absolute, $10^{-7}$ relative, for the whole run. Not the cause. The anisotropic filter now
        takes their floor verbatim anyway, since \S\ref{sec:proposition2} floors
        at ``Mip-Splatting's own value'' and a re-derivation is a needless
        opportunity to differ.
  \item \textbf{The covariance path.} The anisotropic filter passed $\Sigma + \Sigma_{\mathrm{filt}}$
        through the rasteriser's precomputed-covariance input while
        Mip-Splatting passes scales and rotations. The two are algebraically the
        same, but they are different code. Where the filter is isotropic on
        every primitive, the anisotropic filter now takes Mip-Splatting's own path ---
        Proposition~\ref{prop:reduction} used as a computational identity in the code,
        and no longer only as an argument on paper. The gap did not move: $1.245$\,dB. Not the
        cause either.
  \item \textbf{The opacity.} What both paths share is
        Equation~\eqref{eq:filter}'s compensation, and it contained
        \texttt{clamp\_min(1e-12)} on both determinants. $|\Sigma|$ is
        $\prod s_i^2$: a primitive with $s = 0.005$ has $|\Sigma| = 1.6\times
        10^{-14}$, and one with $s=0.001$ has $10^{-18}$. Both determinants
        clamped to the same floor, their ratio became exactly $1$, and the
        compensation switched itself off for most of the model. The anisotropic filter was rendering
        with uncompensated opacity while every baseline number in this thesis
        was produced with it --- which is precisely the observed signature:
        denser and sharper at full resolution, worse at every reduced scale.
\end{enumerate}

With the clamp reduced to a division guard, $|\Sigma|$ taken from the scale
vector where $\prod s_i^2$ is exact, and the isotropic regime using
Mip-Splatting's own scalar expression, the largest difference over the two scenes
then measured fell from $1.245$\,dB to $0.136$\,dB. Over the
$\methodParityScenes$ scenes for which a \emph{paired} measurement exists it is
now $\mathbf{\methodParityMax}$\,dB, comfortably below the seed spread of
\S\ref{sec:spread}.

\keybox{what ``paired'' has to mean here}{%
Proposition~\ref{prop:reduction} says the anisotropic filter \emph{becomes} Mip-Splatting when the
floor dominates, so the two numbers being subtracted have to come from the same
capture, the same harness invocation and the same build; otherwise the
difference also contains whatever separates two runs. That is a stricter
requirement than it sounds, and this table did not meet it. The selector behind
the macro accepted the standard single-scale rows as a stand-in for
Mip-Splatting at the full protocol, so the anisotropic filter was being differenced against a sweep
from three commits earlier that had never used the protocol, with two scenes
entering twice on one side and once on the other. It reported $0.021$\,dB. The
figure was not wrong by much, and that is the point: a comparison can be
mis-specified and still land on a plausible number, which is why the pairing is
now enforced in the generator, and why the scene count is a macro that reads
from the record.}

\keybox{what made this findable}{%
The unit test of \S\ref{sec:proposition2} passed throughout, because its scales
were $0.01$--$0.06$, where $\prod s_i^2$ sits at the clamp and the fault is
nearly invisible. A test whose inputs do not span the real range will certify a
broken component. It now covers $|\Sigma|$ down to $1.4\times10^{-18}$ and prints
that range in its own output. The bug itself was silent --- no error, plausible
numbers, the wrong shape --- and what exposed it was requiring parity to hold \emph{exactly}.}
